\documentclass{article}
\usepackage[margin=1in]{geometry}
\usepackage{amsmath}
\usepackage{amssymb}
\usepackage{bm}
\usepackage{amsthm}

\newtheorem{theorem}{Theorem}
\newtheorem{lemma}{Lemma}
\newtheorem{corollary}{Corollary}
\newtheorem{assumption}{Assumption}
\newtheorem{remark}{Remark}

\begin{document}

\title{Nonparametric Predictive Regression with Exogenous Shocks: A Sieve Empirical Likelihood Approach (In progress)}
\author{Siyan Dong}
\date{\today}
\maketitle

\section{Introduction and Data-Generating Process}

We consider the problem of testing return predictability in the presence of
highly persistent predictors and stationary exogenous covariates. The
exogenous covariate may be a scalar or $d$-dimensional vector of
macroeconomic, environmental, or other predetermined state variables whose
effect on returns is not well described by a constant intercept. The
data-generating process (DGP) is specified by the
following system of equations:

\begin{enumerate}
    \item \textbf{Predictive Regression:}
    \begin{equation}
        Y_t = m(W_{t-1}) + \beta X_{t-1} + U_t
    \end{equation}
    \begin{equation}
        X_t = \theta + \rho X_{t-1} + \varepsilon_t
    \end{equation}
    
    \item \textbf{Endogeneity:}
    \begin{equation}
        U_t = \phi \varepsilon_t + z_t \quad \text{with } E[z_t \mid \varepsilon_t] = 0
    \end{equation}
\end{enumerate}
where $Y_t$ represents asset returns, $W_{t-1} \in \mathbb{R}^d$ is a
stationary exogenous covariate vector, and $X_{t-1}$ is a persistent financial
predictor. The function
$m(W_{t-1})$ is an unknown nonparametric nuisance component: it can be viewed
as a flexible intercept or exogenous-control term that absorbs the direct
effect of $W_{t-1}$ on returns. The key timing restriction is contemporaneous
sieve orthogonality between the exogenous covariate vector and the persistent
predictor weight at the sampling frequency used for inference. This
restriction does not rule out effects of $W_{t-1}$ on returns through
$m(W_{t-1})$, nor does it rule out adjustment in future values of $X_t$.

Throughout the theoretical development, $W_{t-1}$ should be read as a generic
stationary exogenous covariate vector and is not restricted to be univariate;
the scalar case corresponds to $d=1$. The weather-shock interpretation used
below is one empirical application of this more general setup.

Motivated by Perron (1989), our goal is to investigate return predictability while avoiding size distortions or false inferences caused by near unit roots in the predictor dynamics (i.e., $\rho \approx 1$). This is achieved by testing the null hypothesis of no predictability:
\begin{equation} 
\begin{aligned}
    H_0: \quad &\beta = 0 \\
    H_1: \quad &\beta \neq 0
\end{aligned}
\end{equation}

We propose a sieve-orthogonalized empirical likelihood (EL) framework that achieves a standard $\chi^2_1$ limiting distribution under the null under high-level persistent-predictor score conditions.

\section{Empirical Interpretation: Weather Shocks and Cocoa Futures}

The framework can be applied to agricultural commodity markets by interpreting
$W_{t-1}$ as one or more local weather variables, such as precipitation,
temperature, or humidity, and $Y_t$ as excess returns on cocoa futures. The 2023--2024
cocoa market provides a useful example: severe weather conditions in West
Africa contributed to supply deficits and large price movements, while lagged
prices and related financial predictors remained highly persistent.

This empirical setting fits the general framework for three reasons. First,
weather variables are plausibly exogenous to contemporaneous financial-market
innovations at the daily sampling frequency, while still being allowed to
affect subsequent returns through the structural component $m(W_{t-1})$.
Second, the relation between weather and supply is nonlinear, so a flexible
function $m(W_{t-1})$ is more appropriate than a constant intercept or a
linear control. Third, financial predictors such as lagged prices can be
highly persistent, so inference on $\beta$ must remain valid across
near-unit-root persistence regimes.

Thus, in the cocoa application, the generic nuisance component $m(W_{t-1})$
has the interpretation of a nonlinear weather-response surface. The econometric
method itself, however, does not depend on weather specifically; it only
requires a stationary exogenous covariate vector satisfying the orthogonality
and regularity conditions stated below.

\section{The Orthogonalized Sieve-EL Framework}

We handle the unknown nonparametric nuisance component $m(W_{t-1})$ by approximating it using a sieve basis of dimension $K$ (e.g., polynomials or splines). Let $\bm{P}_K(W_{t-1}) = (p_1(W_{t-1}), \dots, p_K(W_{t-1}))'$ be the $K \times 1$ vector of these basis functions; when $d>1$, the $p_j$ are multivariate basis functions.

\subsection{The Step-by-Step Algorithm}

The estimation and inference procedure is implemented sequentially via the following three steps. The key idea is to profile out the flexible intercept component and then construct an EL moment that is orthogonal to the same sieve space.

\subsubsection*{Step 1: Profiling out the Nonparametric Nuisance Component}
For each candidate parameter $\beta$ under the null hypothesis, we profile out the unknown function $m(W_{t-1})$, which plays the role of a flexible intercept or exogenous-control component. Let $\bm{Y} = (Y_1, \dots, Y_T)'$ and $\bm{X}_{lag} = (X_0, \dots, X_{T-1})'$. Let $\bm{P}$ be the $T \times K$ matrix whose $t$-th row is $\bm{P}_K(W_{t-1})'$. The profile sieve coefficients for the nuisance component are:
\begin{equation}
    \widehat{\bm{c}}_K(\beta) = (\bm{P}'\bm{P})^{-1}\bm{P}'(\bm{Y} - \beta \bm{X}_{lag})
\end{equation}
The profile residual is defined using the sieve annihilator matrix $\bm{M}_K = \bm{I}_T - \bm{P}(\bm{P}'\bm{P})^{-1}\bm{P}'$:
\begin{equation}
    \widehat{\bm{u}}(\beta) = \bm{M}_K(\bm{Y} - \beta \bm{X}_{lag})
\end{equation}
Under the null hypothesis $\beta = \beta_0$, substituting the true DGP $\bm{Y} = \bm{P}\bm{c} + \bm{r} + \beta_0 \bm{X}_{lag} + \bm{U}$ (where $\bm{r}$ is the approximation error) yields:
\begin{equation}
    \widehat{\bm{u}}(\beta_0) = \bm{M}_K(\bm{P}\bm{c} + \bm{r} + \bm{U}) = \bm{M}_K(\bm{r} + \bm{U})
\end{equation}

\subsubsection*{Step 2: The Sieve Weight Projection}
In standard Empirical Likelihood, estimating a high-dimensional nuisance component can affect the limiting distribution. To neutralize this generated-nuisance effect, we construct an orthogonalized weight function. Let $\bm{w} = (w_1, \dots, w_T)'$ where $w_t = w(X_{t-1})$ is the raw weight (e.g., $\tanh(X_{t-1}/c_w)$). We project the weights onto the same sieve basis and define the orthogonalized weight $\bm{w}^c$ as the residual:
\begin{equation}
    \bm{w}^c = \bm{M}_K \bm{w}
\end{equation}
Let $w_t^c = (\bm{w}^c)_t$. By the properties of orthogonal projection, the orthogonalized weights are exactly orthogonal to the sieve basis in-sample:
\begin{equation}
    \sum_{t=1}^T \bm{P}_K(W_{t-1}) w_t^c \equiv \bm{0}
\end{equation}

\subsubsection*{Step 3: Algebraic Elimination in the EL Moment Condition}
For testing $H_0: \beta = \beta_0$, the EL moment condition is evaluated as:
\begin{equation}
    \widehat{Z}_t(\beta_0) = \widehat{u}_t(\beta_0) w_t^c
\end{equation}
Summing the moment conditions across the sample to evaluate the numerator of the EL statistic:
\begin{equation}
    \sum_{t=1}^T \widehat{Z}_t(\beta_0) = \widehat{\bm{u}}(\beta_0)' \bm{w}^c = (\bm{r} + \bm{U})' \bm{M}_K \bm{M}_K \bm{w} = (\bm{r} + \bm{U})' \bm{w}^c
\end{equation}
The fitted sieve component is thereby eliminated algebraically. The remaining approximation residual is controlled by the sieve rate conditions, leaving:
\begin{equation}
    \frac{1}{\sqrt{T}} \sum_{t=1}^T \widehat{Z}_t(\beta_0) = \frac{1}{\sqrt{T}} \sum_{t=1}^T U_t w_t^c + \underbrace{\frac{1}{\sqrt{T}} \sum_{t=1}^T r_t w_t^c}_{=\; o_p(1)}
\end{equation}
Because the moment condition is orthogonal to the fitted nuisance space, the sample variance of $\widehat{Z}_t(\beta_0)$ consistently estimates the relevant innovation variance under the stated rate conditions. As a result, the Empirical Likelihood ratio statistic (where $\widehat{\lambda}$ is the Lagrange multiplier) achieves the standard chi-squared limit:
\begin{equation}
    \ell_T(\beta_0) = 2 \sum_{t=1}^T \log(1 + \widehat{\lambda} \widehat{Z}_t(\beta_0)) \xrightarrow{d} \chi^2_1
\end{equation}

\subsection{Relation to Fixed-Intercept Predictive EL}

The framework differs from fixed-intercept predictive EL in two fundamental aspects:
\begin{itemize}
    \item \textbf{Nuisance Dimensionality}: A fixed scalar intercept $\alpha$ is replaced by an infinite-dimensional nuisance function $m(W_{t-1})$, which requires a sieve approximation of dimension $K \to \infty$.
    \item \textbf{Orthogonalization Space}: The demeaning against a constant $1$ is replaced by residualizing against the sieve basis $\bm{P}_K(W)$. The orthogonalized weight $\bm{w}^c$ is projected onto the $K$-dimensional sieve space $\text{span}\{\bm{P}_K(W_{t-1})\}$. We keep the AR(1) equation as the persistence DGP, but we do not use AR(1) OLS residuals in the EL statistic.
\end{itemize}

\section{The Growth Rate of the Sieve Dimension $K$}

To establish the asymptotic validity of our test statistic, we must place bounds on the growth rate of the tuning parameter $K$ relative to the sample size $T$.

\subsection{Lower Bound (Bias Elimination)}
Assume the true nuisance function $m(W)$ has smoothness $s$ and that $W$ is $d$-dimensional. Sieve theory implies that the approximation error $r_t$ satisfies $\sup |r_t| = O_p(K^{-s/d})$. To ensure that the approximation bias does not distort the asymptotic distribution, we require $T^{-1/2} \sum_{t=1}^T r_t w^c_t = o_p(1)$, which translates to:
\begin{equation}
    \sqrt{T} K^{-s/d} \to 0 \quad \implies \quad K \gg T^{\frac{d}{2s}}
\end{equation}

\subsection{Upper Bound (Variance Control and Leakage Mitigation)}
Estimating $K$ nuisance parameters introduces in-sample variance inflation and potential projection leakage. Under mixing conditions for the exogenous covariate process $W_t$, this leakage does not accumulate, and the leading variance penalty scales as $K / \sqrt{T}$. To prevent variance inflation from affecting the EL ratio, we require:
\begin{equation}
    \frac{K}{\sqrt{T}} \to 0 \quad \implies \quad K \ll T^{\frac{1}{2}}
\end{equation}

\subsection{Growth rate of K}
Combining these bounds yields the formal asymptotic growth condition:
\begin{equation}
    T^{\frac{d}{2s}} \ll K \ll T^{\frac{1}{2}}
\end{equation}
This rate window is nonempty when $s > d$.

% \section{Formal Orthogonality Assumptions}

% To ensure the algebraic orthogonalization rigorously eliminates the fitted
% nonparametric nuisance component, we must state the relationship between the
% persistent financial predictor and the stationary exogenous covariate vector. We
% impose the following two conditions on $W_{t-1}$:

% \paragraph*{Assumption 5.1 (Strict Exogeneity to the Market Innovation)} \textit{The exogenous covariate vector $W_{t-1}$ is strictly exogenous to the financial market innovation $U_t$. Mathematically:}
% \begin{equation}
%     \mathbb{E}[U_t \mid \mathcal{F}_{t-1}, W_{t-1}] = 0.
% \end{equation}

% \noindent \textbf{Remark.} This condition says the exogenous covariate vector is not a proxy for the residual market innovation $U_t$ after conditioning on predetermined market information. It still allows $W_{t-1}$ to affect returns through the structural component $m(W_{t-1})$ and to be incorporated into future predictor dynamics.

% \paragraph*{Assumption 5.2 (Contemporaneous Sieve Orthogonality)} \textit{At the sampling frequency used for inference, the exogenous covariate vector is contemporaneously orthogonal to the predetermined persistent predictor in the sieve directions used for profiling. For the EL weight $w(X_{t-1})$,}
% \begin{equation}
%     \mathbb{E}[w(X_{t-1}) \mid W_{t-1}] = \mathbb{E}[w(X_{t-1})].
% \end{equation}
% \textit{Equivalently, for any centered nonconstant sieve vector $\bm R_K(W_{t-1})$,}
% \begin{equation}
%     \mathbb{E}\left[\bm R_K(W_{t-1})\{w(X_{t-1})-\mathbb{E}w(X_{t-1})\}\right]=\bm0.
% \end{equation}
% \textit{This is a contemporaneous timing restriction, not a dynamic exclusion: $W_{t-1}$ may affect $Y_t$ through $m(W_{t-1})$ and may affect subsequent returns and future predictor values.}

% \noindent \textbf{Remark.} This assumption is the population version of the sample orthogonality condition used in the proof, $\|T^{-1}\bm R'\widetilde{\bm w}\|=O_p(\sqrt{K/T})$. It preserves the persistent predictor signal at the moment used for inference while allowing the exogenous covariate to enter later market adjustment.

% \section{Wilks' Theorem for Mean Predictability}
% \label{sec:mean-wilks}

% This section gives a formal Wilks theorem for the mean-predictability test.  It is
% useful to write the nuisance step in null-imposed profile form.  Let
% \[
%     \bm Y=(Y_1,\ldots,Y_T)',\qquad
%     \bm X_{lag}=(X_0,\ldots,X_{T-1})',
% \]
% and let \(\bm P\) be the \(T\times K\) matrix whose \(t\)-th row is
% \(\bm P_K(W_{t-1})'\).  Define
% \[
%     \bm\Pi_K=\bm P(\bm P'\bm P)^{-1}\bm P',
%     \qquad
%     \bm M_K=\bm I_T-\bm\Pi_K.
% \]
% For each candidate value \(\beta\), define the null-imposed profile sieve
% coefficient and the corresponding residual by
% \begin{equation}
%     \widehat{\bm c}_K(\beta)
%     =(\bm P'\bm P)^{-1}\bm P'(\bm Y-\beta\bm X_{lag}),
%     \qquad
%     \widehat{\bm u}(\beta)
%     =\bm M_K(\bm Y-\beta\bm X_{lag}).
%     \label{eq:profile-sieve}
% \end{equation}
% Let \(w_t=w(X_{t-1})\), \(\bm w=(w_1,\ldots,w_T)'\), and define the
% sieve-orthogonalized weight
% \begin{equation}
%     \bm w^c=\bm M_K\bm w,
%     \qquad
%     w_t^c=(\bm w^c)_t.
%     \label{eq:sieve-weight}
% \end{equation}
% The scalar estimating equation is
% \begin{equation}
%     \widehat Z_t(\beta)
%     =\widehat u_t(\beta)w_t^c.
%     \label{eq:profile-score}
% \end{equation}
% The empirical likelihood and its log-likelihood ratio are
% \begin{align}
%     L_T(\beta)
%     &=\sup\left\{
%        \prod_{t=1}^T(T\pi_t):
%        \pi_t\geq0,\quad
%        \sum_{t=1}^T\pi_t=1,\quad
%        \sum_{t=1}^T\pi_t\widehat Z_t(\beta)=0
%        \right\},                                                    \label{eq:profile-el}\\
%     \ell_T(\beta)
%     &=-2\log L_T(\beta)
%       =2\sum_{t=1}^T\log\{1+\widehat\lambda(\beta)
%           \widehat Z_t(\beta)\},                                    \label{eq:profile-elr}
% \end{align}
% where \(\widehat\lambda(\beta)\) solves
% \begin{equation}
%     \sum_{t=1}^T
%        \frac{\widehat Z_t(\beta)}
%             {1+\widehat\lambda(\beta)\widehat Z_t(\beta)}=0.
%     \label{eq:lambda-equation}
% \end{equation}

% For a vector \(\bm a=(a_1,\ldots,a_T)'\), write
% \(\|\bm a\|_T^2=T^{-1}\sum_{t=1}^T a_t^2\).  Also write
% \[
%     \overline w_T=T^{-1}\sum_{t=1}^T w_t,
%     \qquad
%     \widetilde w_t=w_t-\overline w_T.
% \]

% \begin{assumption}[Persistent-predictor score conditions]
% \label{ass:score-core}
% The predictor satisfies
% \(X_t=\theta+\rho_T X_{t-1}+\varepsilon_t\), where \(\rho_T\) belongs to
% one of the admissible persistence classes
% including the stationary, mildly integrated, local-to-unity,
% integrated, and mildly explosive cases.  With respect to the filtration
% \[
%     \mathcal H_t
%     =\sigma\{(z_s,\varepsilon_s):s\leq t\}
%        \vee\sigma\{W_s:s\leq t\},
% \]
% \(E(U_t\mid\mathcal H_{t-1})=0\),
% \(\sup_tE|U_t|^{2+\delta}<\infty\) for some \(\delta>0\), and the
% conditional-variance and weight conditions needed for the oracle score limit
% hold.  In particular, for
% \[
%     A_T=\frac1{\sqrt T}\sum_{t=1}^T U_t\widetilde w_t,
%     \qquad
%     V_T=\frac1T\sum_{t=1}^T U_t^2\widetilde w_t^2,
% \]
% there exists an a.s. positive, possibly random, variable \(\eta^2\) such that
% \begin{equation}
%     A_T\ \xrightarrow{\mathrm{st}}\ MN(0,\eta^2),
%     \qquad
%     V_T\ \xrightarrow{p}\ \eta^2,
%     \qquad
%     \max_{1\leq t\leq T}|U_t\widetilde w_t|=o_p(\sqrt T).
%     \label{eq:score-core}
% \end{equation}
% The raw weight is uniformly bounded and satisfies the stated saturation
% condition.
% \end{assumption}

% \begin{assumption}[Exogenous-covariate sieve orthogonality and projection]
% \label{ass:sieve-projection}
% The first component of \(\bm P_K(W)\) is one.  After centering the remaining
% basis functions in sample, write the corresponding \(T\times(K-1)\) matrix
% as \(\bm R\), so that \(\bm 1_T'\bm R=\bm0\) and
% \(\operatorname{span}(\bm P)=\operatorname{span}(\bm1_T,\bm R)\).  The
% primitive timing restriction is contemporaneous sieve orthogonality:
% \[
%     \mathbb E\left[\bm R_K(W_{t-1})
%     \{w(X_{t-1})-\mathbb E w(X_{t-1})\}\right]=\bm0,
% \]
% where \(\bm R_K(W_{t-1})\) denotes the population counterpart of the centered
% nonconstant sieve vector.  This condition allows \(W_{t-1}\) to affect
% \(Y_t\) and future values of \(X\); it only rules out contemporaneous
% projection of the predictor weight onto the exogenous-covariate sieve.  The
% strict-exogeneity condition in Assumption~\ref{ass:score-core} holds with
% contemporaneous exogenous-covariate information included in the conditioning set.

% There are constants \(0<c<C<\infty\) such that, with probability tending to
% one,
% \[
%     c\leq\lambda_{\min}(T^{-1}\bm R'\bm R)
%       \leq\lambda_{\max}(T^{-1}\bm R'\bm R)\leq C.
% \]
% Let \(\zeta_K=\max_{t\leq T}\|\bm R_t\|\), where \(\bm R_t'\) is row \(t\)
% of \(\bm R\).  The sample analogue of the population orthogonality and the
% remaining increasing-dimensional score bounds are:
% \begin{align}
%     \left\|T^{-1}\bm R'\widetilde{\bm w}\right\|
%         &=O_p\!\left(\sqrt{K/T}\right),                                  \label{eq:rate-Rw}\\
%     \left\|T^{-1/2}\bm R'\bm U\right\|&=O_p(\sqrt K),                    \label{eq:rate-RU}\\
%     \left\|T^{-1}\bm P'\bm U\right\|
%         &=O_p\!\left(\sqrt{K/T}\right),                                  \label{eq:rate-PU}
% \end{align}
% where \(\widetilde{\bm w}=(\widetilde w_1,\ldots,\widetilde w_T)'\) and
% \(\bm U=(U_1,\ldots,U_T)'\).  These bounds follow under the stated sieve
% orthogonality, exponential mixing of \(W_t\), uniformly bounded basis moments,
% and the martingale moment condition in Assumption~\ref{ass:score-core}.
% \end{assumption}

% \begin{assumption}[Approximation and growth rates]
% \label{ass:approximation}
% For some \(\bm c_K\),
% \[
%     m(W_{t-1})=\bm P_K(W_{t-1})'\bm c_K+r_{K,t},
% \]
% where
% \[
%     \|\bm r_K\|_T=O(a_K),\qquad
%     \|\bm r_K\|_\infty=O(a_K),\qquad
%     \sqrt T\,a_K\longrightarrow0.
% \]
% The sieve projector is stable in sup norm on the approximation residual, so
% that \(\|\bm M_K\bm r_K\|_\infty=O(a_K)\).  Moreover,
% \begin{equation}
%     K/\sqrt T\longrightarrow0,
%     \qquad
%     \zeta_K\sqrt{K/T}\longrightarrow0.
%     \label{eq:K-rates}
% \end{equation}
% For normalized spline bases, for which typically \(\zeta_K\asymp\sqrt K\),
% the second condition in \eqref{eq:K-rates} is implied by
% \(K/\sqrt T\to0\).  If \(a_K\asymp K^{-s/d}\), a sufficient rate window is
% \(T^{d/(2s)}\ll K\ll T^{1/2}\), which is nonempty when \(s>d\).
% \end{assumption}

% \begin{assumption}[Nondegeneracy and convex hull]
% \label{ass:convex-hull}
% There exists \(\underline\eta>0\) such that
% \(P(\eta^2\geq\underline\eta^2)=1\).  In addition,
% \[
%     P\left\{
%       \min_{1\leq t\leq T}\widehat Z_t(\beta_0)<0<
%       \max_{1\leq t\leq T}\widehat Z_t(\beta_0)
%     \right\}\longrightarrow1.
% \]
% \end{assumption}

% \begin{lemma}[Negligibility of the growing sieve projection]
% \label{lem:projection-negligible}
% Under Assumptions~\ref{ass:sieve-projection}--\ref{ass:approximation},
% \begin{align}
%     \|\bm w^c-\widetilde{\bm w}\|_T
%        &=O_p\!\left(\sqrt{K/T}\right),                                   \label{eq:weight-L2}\\
%     \max_{t\leq T}|w_t^c-\widetilde w_t|
%        &=O_p\!\left(\zeta_K\sqrt{K/T}\right)=o_p(1),                    \label{eq:weight-max}\\
%     \frac1{\sqrt T}\sum_{t=1}^T
%        U_t(w_t^c-\widetilde w_t)
%        &=O_p(K/\sqrt T)=o_p(1),                                           \label{eq:weight-score}\\
%     \|\bm\Pi_K\bm U\|_T
%        &=O_p\!\left(\sqrt{K/T}\right),                                  \label{eq:U-proj-L2}\\
%     \max_{t\leq T}|(\bm\Pi_K\bm U)_t|
%        &=O_p\!\left(\zeta_K\sqrt{K/T}\right)=o_p(1).                   \label{eq:U-proj-max}
% \end{align}
% \end{lemma}

% \begin{proof}
% Because the sieve contains a constant and \(\bm1_T'\bm R=\bm0\), residualizing
% \(\bm w\) on \(\bm P\) is equivalent to residualizing
% \(\widetilde{\bm w}\) on \(\bm R\).  Hence
% \[
%     \bm w^c=\widetilde{\bm w}-\bm R\widehat{\bm\delta}_K,
%     \qquad
%     \widehat{\bm\delta}_K
%        =(\bm R'\bm R)^{-1}\bm R'\widetilde{\bm w}.
% \]
% The Gram-matrix condition and \eqref{eq:rate-Rw} imply
% \(\|\widehat{\bm\delta}_K\|=O_p(\sqrt{K/T})\).  Therefore,
% \[
%     \|\bm R\widehat{\bm\delta}_K\|_T^2
%        =\widehat{\bm\delta}_K'
%           (T^{-1}\bm R'\bm R)\widehat{\bm\delta}_K
%        =O_p(K/T),
% \]
% which proves \eqref{eq:weight-L2}; and
% \[
%     \max_t|\bm R_t'\widehat{\bm\delta}_K|
%        \leq\zeta_K\|\widehat{\bm\delta}_K\|
%        =O_p\!\left(\zeta_K\sqrt{K/T}\right),
% \]
% which proves \eqref{eq:weight-max}.  Moreover,
% \[
%     \frac1{\sqrt T}\bm U'(\bm w^c-\widetilde{\bm w})
%       =-\left(T^{-1/2}\bm R'\bm U\right)'
%           \widehat{\bm\delta}_K
%       =O_p(\sqrt K)O_p\!\left(\sqrt{K/T}\right)
%       =O_p(K/\sqrt T),
% \]
% giving \eqref{eq:weight-score}.

% Finally, let
% \(\widehat{\bm a}_U=(\bm P'\bm P)^{-1}\bm P'\bm U\).  The Gram condition
% and \eqref{eq:rate-PU} imply
% \(\|\widehat{\bm a}_U\|=O_p(\sqrt{K/T})\).  The same quadratic-form and
% maximum-row arguments give \eqref{eq:U-proj-L2} and
% \eqref{eq:U-proj-max}.
% \end{proof}

% \begin{lemma}[Asymptotic equivalence of the feasible and oracle scores]
% \label{lem:score-equivalence}
% Under Assumptions~\ref{ass:score-core}--\ref{ass:approximation}, under
% \(H_0:\beta=\beta_0\),
% \begin{align}
%     S_T
%       :=\frac1{\sqrt T}\sum_{t=1}^T\widehat Z_t(\beta_0)
%       &=\frac1{\sqrt T}\sum_{t=1}^T U_t\widetilde w_t+o_p(1),             \label{eq:numerator-equivalence}\\
%     \widehat V_T
%       :=\frac1T\sum_{t=1}^T\widehat Z_t(\beta_0)^2
%       &=\frac1T\sum_{t=1}^T U_t^2\widetilde w_t^2+o_p(1),                \label{eq:variance-equivalence}\\
%     \max_{t\leq T}|\widehat Z_t(\beta_0)|&=o_p(\sqrt T).                \label{eq:max-score}
% \end{align}
% Consequently, \(\widehat V_T\xrightarrow{p}\eta^2\).
% \end{lemma}

% \begin{proof}
% Under the null,
% \[
%     \bm Y-\beta_0\bm X_{lag}
%        =\bm P\bm c_K+\bm r_K+\bm U,
%     \qquad
%     \widehat{\bm u}(\beta_0)
%        =\bm M_K(\bm U+\bm r_K).
% \]
% Because \(\bm M_K\) is symmetric and idempotent and
% \(\bm w^c=\bm M_K\bm w\),
% \begin{align*}
%     S_T
%       &=T^{-1/2}\widehat{\bm u}(\beta_0)'\bm w^c\\
%       &=T^{-1/2}(\bm U+\bm r_K)'\bm M_K^2\bm w\\
%       &=T^{-1/2}(\bm U+\bm r_K)'\bm w^c.
% \end{align*}
% Lemma~\ref{lem:projection-negligible} gives
% \(T^{-1/2}\bm U'\bm w^c
%   =T^{-1/2}\bm U'\widetilde{\bm w}+o_p(1)\).  Also, by Cauchy--Schwarz and
% the contraction property of an orthogonal projection,
% \[
%     \left|T^{-1/2}\bm r_K'\bm w^c\right|
%        \leq\sqrt T\,\|\bm r_K\|_T\|\bm w^c\|_T
%        \leq\sqrt T\,\|\bm r_K\|_T\|\bm w\|_T
%        =o_p(1).
% \]
% This proves \eqref{eq:numerator-equivalence}.

% For the denominator result, define the feasible-oracle residuals
% \[
%     \Delta_{u,t}=\widehat u_t(\beta_0)-U_t,
%     \qquad
%     \Delta_{w,t}=w_t^c-\widetilde w_t,
%     \qquad
%     \Delta_{Z,t}=\widehat Z_t(\beta_0)-U_t\widetilde w_t.
% \]
% Then
% \[
%     \bm\Delta_u=-\bm\Pi_K\bm U+\bm M_K\bm r_K,
% \]
% so that, by Lemma~\ref{lem:projection-negligible} and projection contraction,
% \[
%     \|\bm\Delta_u\|_T
%        =O_p\!\left(\sqrt{K/T}+a_K\right)=o_p(1).
% \]
% Moreover, \(\max_t|\Delta_{w,t}|=o_p(1)\),
% \(\max_t|w_t^c|=O_p(1)\), and \(T^{-1}\sum_tU_t^2=O_p(1)\).  Since
% \[
%     \Delta_{Z,t}
%        =\Delta_{u,t}w_t^c+U_t\Delta_{w,t},
% \]
% it follows that
% \begin{align*}
%     \|\bm\Delta_Z\|_T^2
%     &\leq
%        2\max_t|w_t^c|^2\|\bm\Delta_u\|_T^2
%        +2\max_t|\Delta_{w,t}|^2\frac1T\sum_{t=1}^TU_t^2\\
%     &=o_p(1).
% \end{align*}
% Let \(V_T=T^{-1}\sum_{t=1}^TU_t^2\widetilde w_t^2\).  Then
% \[
%     \widehat V_T-V_T
%       =\frac2T\sum_{t=1}^T U_t\widetilde w_t\Delta_{Z,t}
%         +\|\bm\Delta_Z\|_T^2.
% \]
% By Cauchy--Schwarz,
% \[
%     |\widehat V_T-V_T|
%       \leq 2V_T^{1/2}\|\bm\Delta_Z\|_T+\|\bm\Delta_Z\|_T^2
%       =o_p(1),
% \]
% because \(V_T=O_p(1)\) by Assumption~\ref{ass:score-core}.  This proves
% \eqref{eq:variance-equivalence}; the displayed convergence of
% \(\widehat V_T\) follows from \(V_T\xrightarrow{p}\eta^2\) in
% Assumption~\ref{ass:score-core}.

% Finally, \(\max_t|U_t|=o_p(\sqrt T)\) follows from the uniform
% \((2+\delta)\)-moment bound.  Together with
% \eqref{eq:U-proj-max}, the sup-norm stability of \(\bm M_K\bm r_K\), and
% \(\max_t|w_t^c|=O_p(1)\), this proves \eqref{eq:max-score}.
% \end{proof}

% \begin{theorem}[Wilks phenomenon for sieve-orthogonalized mean EL]
% \label{thm:sieve-wilks}
% Suppose Assumptions~\ref{ass:score-core}--\ref{ass:convex-hull} hold.  Then,
% under \(H_0:\beta=\beta_0\),
% \begin{equation}
%     \ell_T(\beta_0)\ \xrightarrow{d}\ \chi_1^2.
%     \label{eq:wilks-result}
% \end{equation}
% The conclusion holds for each predictor-persistence class covered by the high-level
% mean theorem.
% \end{theorem}

% \begin{proof}
% By Lemma~\ref{lem:score-equivalence} and
% Assumption~\ref{ass:score-core},
% \[
%     S_T\ \xrightarrow{\mathrm{st}}\ MN(0,\eta^2),
%     \qquad
%     \widehat V_T\ \xrightarrow{p}\eta^2.
% \]
% Consequently,
% \begin{equation}
%     \frac{S_T}{\sqrt{\widehat V_T}}
%        \ \xrightarrow{\mathrm{st}}\ N(0,1),
%     \qquad
%     \frac{S_T^2}{\widehat V_T}
%        \ \xrightarrow{d}\chi_1^2.
%     \label{eq:self-normalized-limit}
% \end{equation}

% It remains to connect the self-normalized statistic to empirical likelihood.
% Write \(\widehat Z_t=\widehat Z_t(\beta_0)\) and
% \(\widehat\lambda=\widehat\lambda(\beta_0)\).  The convex-hull condition,
% \(S_T=O_p(1)\), \(\widehat V_T\to_p\eta^2>0\), and
% \(\max_t|\widehat Z_t|=o_p(\sqrt T)\) imply the standard scalar-EL bounds
% \[
%     \widehat\lambda=O_p(T^{-1/2}),
%     \qquad
%     \max_{t\leq T}|\widehat\lambda\widehat Z_t|=o_p(1).
% \]
% Expanding the multiplier equation \eqref{eq:lambda-equation} gives
% \[
%     0=\sum_{t=1}^T\widehat Z_t
%       -\widehat\lambda\sum_{t=1}^T\widehat Z_t^2
%       +o_p(\sqrt T),
% \]
% and hence
% \begin{equation}
%     \widehat\lambda
%        =\frac{\sum_{t=1}^T\widehat Z_t}
%               {\sum_{t=1}^T\widehat Z_t^2}
%         +o_p(T^{-1/2}).
%     \label{eq:lambda-expansion}
% \end{equation}
% A second-order expansion of \eqref{eq:profile-elr}, together with
% \(\max_t|\widehat\lambda\widehat Z_t|=o_p(1)\), yields
% \begin{align*}
%     \ell_T(\beta_0)
%       &=2\widehat\lambda\sum_{t=1}^T\widehat Z_t
%         -\widehat\lambda^2\sum_{t=1}^T\widehat Z_t^2+o_p(1)\\
%       &=\frac{\left(T^{-1/2}\sum_{t=1}^T\widehat Z_t\right)^2}
%               {T^{-1}\sum_{t=1}^T\widehat Z_t^2}
%         +o_p(1)\\
%       &=\frac{S_T^2}{\widehat V_T}+o_p(1).
% \end{align*}
% The result follows from \eqref{eq:self-normalized-limit} and Slutsky's
% theorem.
% \end{proof}

% \begin{corollary}[Test of no mean predictability]
% \label{cor:mean-test}
% For \(H_0:\beta=0\), reject at asymptotic level \(\alpha\) whenever
% \[
%     \ell_T(0)>\chi^2_{1,1-\alpha}.
% \]
% Equivalently, the asymptotic \(p\)-value is
% \(1-F_{\chi_1^2}\{\ell_T(0)\}\).  A confidence set for \(\beta\) is obtained
% by inversion:
% \[
%     \mathcal C_{1-\alpha}
%       =\{\beta:\ell_T(\beta)\leq\chi^2_{1,1-\alpha}\}.
% \]
% \end{corollary}

% \begin{remark}[Relation to the unrestricted two-stage estimator]
% \label{rem:unrestricted-chat}
% The profile construction \eqref{eq:profile-sieve} is the cleanest route to
% Theorem~\ref{thm:sieve-wilks}: at \(\beta_0\), the fitted nuisance error is
% exactly \(\bm\Pi_K(\bm U+\bm r_K)\), whose prediction norm is
% \(O_p(\sqrt{K/T}+a_K)\).  If one instead retains a separately estimated
% unrestricted coefficient \(\widehat{\bm c}\), the same theorem continues to
% hold provided one proves the additional generated-nuisance condition
% \[
%     \frac1T\sum_{t=1}^T
%        \left\{\bm P_K(W_{t-1})'
%               (\widehat{\bm c}-\bm c_K)\right\}^2(w_t^c)^2=o_p(1)
% \]
% and the corresponding maximum-term bound.  Those statements do not follow
% from the algebraic identity \(\bm P'\bm w^c=\bm0\) alone.
% \end{remark}

% \begin{remark}[Why the same argument is not immediately a quantile theorem]
% \label{rem:quantile}
% For a quantile score \(\psi_\tau(u)=\tau-\mathbf1\{u<0\}\), the nuisance error
% enters nonlinearly:
% \[
%   \psi_\tau\!\left(U_{t,\tau}+r_{K,t}
%        -\bm P_K(W_{t-1})'(\widehat{\bm c}_\tau-\bm c_{K,\tau})\right).
% \]
% Therefore \(\bm P'\bm w^c=\bm0\) does not algebraically cancel the first-stage
% error.  A first-order expansion instead produces a density-weighted term of
% the form
% \[
%   \sum_{t=1}^T f_{t,\tau}(0)
%        \bm P_K(W_{t-1})w_t^c,
% \]
% which is not generally zero.  A quantile extension would require either a
% restrictive condition making \(f_{t,\tau}(0)\) constant with respect to the
% predictor, or a new density-weighted orthogonalization and additional
% nonparametric-rate arguments.  Thus Theorem~\ref{thm:sieve-wilks} is a
% complete and natural endpoint for the mean paper, whereas a general quantile
% extension is a separate methodological problem.
% \end{remark}

% \section{Conclusion}

% Standard predictive regressions often absorb nonlinear exogenous state
% variables into a fixed intercept or the regression error. That simplification
% can leave important nuisance variation in the residual and can interact poorly
% with highly persistent predictors. This paper instead treats the intercept
% component as an unknown function $m(W)$, profiles it out under each candidate
% $\beta$, and residualizes the EL weight against the same sieve space.

% The AR(1) specification is used only to describe the persistence class of the
% predictor. The empirical likelihood statistic itself does not estimate or
% include AR residuals. The resulting orthogonalized Sieve-EL statistic provides
% standard inference on mean predictability while allowing a flexible
% exogenous-control component. Weather shocks in commodity markets are one
% motivating application, but the theory applies more generally to stationary
% exogenous covariates satisfying the stated orthogonality and regularity
% conditions.

\end{document}
